% Options for packages loaded elsewhere
% Options for packages loaded elsewhere
\PassOptionsToPackage{unicode}{hyperref}
\PassOptionsToPackage{hyphens}{url}
\PassOptionsToPackage{dvipsnames,svgnames,x11names}{xcolor}
%
\documentclass[
  english,
  11pt,
]{article}
\usepackage{xcolor}
\usepackage[margin=25mm]{geometry}
\usepackage{amsmath,amssymb}
\setcounter{secnumdepth}{-\maxdimen} % remove section numbering
\usepackage{iftex}
\ifPDFTeX
  \usepackage[T1]{fontenc}
  \usepackage[utf8]{inputenc}
  \usepackage{textcomp} % provide euro and other symbols
\else % if luatex or xetex
  \usepackage{unicode-math} % this also loads fontspec
  \defaultfontfeatures{Scale=MatchLowercase}
  \defaultfontfeatures[\rmfamily]{Ligatures=TeX,Scale=1}
\fi
\usepackage{lmodern}
\ifPDFTeX\else
  % xetex/luatex font selection
\fi
% Use upquote if available, for straight quotes in verbatim environments
\IfFileExists{upquote.sty}{\usepackage{upquote}}{}
\IfFileExists{microtype.sty}{% use microtype if available
  \usepackage[]{microtype}
  \UseMicrotypeSet[protrusion]{basicmath} % disable protrusion for tt fonts
}{}
\makeatletter
\@ifundefined{KOMAClassName}{% if non-KOMA class
  \IfFileExists{parskip.sty}{%
    \usepackage{parskip}
  }{% else
    \setlength{\parindent}{0pt}
    \setlength{\parskip}{6pt plus 2pt minus 1pt}}
}{% if KOMA class
  \KOMAoptions{parskip=half}}
\makeatother
% Make \paragraph and \subparagraph free-standing
\makeatletter
\ifx\paragraph\undefined\else
  \let\oldparagraph\paragraph
  \renewcommand{\paragraph}{
    \@ifstar
      \xxxParagraphStar
      \xxxParagraphNoStar
  }
  \newcommand{\xxxParagraphStar}[1]{\oldparagraph*{#1}\mbox{}}
  \newcommand{\xxxParagraphNoStar}[1]{\oldparagraph{#1}\mbox{}}
\fi
\ifx\subparagraph\undefined\else
  \let\oldsubparagraph\subparagraph
  \renewcommand{\subparagraph}{
    \@ifstar
      \xxxSubParagraphStar
      \xxxSubParagraphNoStar
  }
  \newcommand{\xxxSubParagraphStar}[1]{\oldsubparagraph*{#1}\mbox{}}
  \newcommand{\xxxSubParagraphNoStar}[1]{\oldsubparagraph{#1}\mbox{}}
\fi
\makeatother


\usepackage{longtable,booktabs,array}
\newcounter{none} % for unnumbered tables
\usepackage{calc} % for calculating minipage widths
% Correct order of tables after \paragraph or \subparagraph
\usepackage{etoolbox}
\makeatletter
\patchcmd\longtable{\par}{\if@noskipsec\mbox{}\fi\par}{}{}
\makeatother
% Allow footnotes in longtable head/foot
\IfFileExists{footnotehyper.sty}{\usepackage{footnotehyper}}{\usepackage{footnote}}
\makesavenoteenv{longtable}
\usepackage{graphicx}
\makeatletter
\newsavebox\pandoc@box
\newcommand*\pandocbounded[1]{% scales image to fit in text height/width
  \sbox\pandoc@box{#1}%
  \Gscale@div\@tempa{\textheight}{\dimexpr\ht\pandoc@box+\dp\pandoc@box\relax}%
  \Gscale@div\@tempb{\linewidth}{\wd\pandoc@box}%
  \ifdim\@tempb\p@<\@tempa\p@\let\@tempa\@tempb\fi% select the smaller of both
  \ifdim\@tempa\p@<\p@\scalebox{\@tempa}{\usebox\pandoc@box}%
  \else\usebox{\pandoc@box}%
  \fi%
}
% Set default figure placement to htbp
\def\fps@figure{htbp}
\makeatother



\ifLuaTeX
\usepackage[bidi=basic,shorthands=off]{babel}
\else
\usepackage[bidi=default,shorthands=off]{babel}
\fi
\ifLuaTeX
  \usepackage{selnolig} % disable illegal ligatures
\fi


\setlength{\emergencystretch}{3em} % prevent overfull lines

\providecommand{\tightlist}{%
  \setlength{\itemsep}{0pt}\setlength{\parskip}{0pt}}



 


\makeatletter
\@ifpackageloaded{caption}{}{\usepackage{caption}}
\AtBeginDocument{%
\ifdefined\contentsname
  \renewcommand*\contentsname{Table of contents}
\else
  \newcommand\contentsname{Table of contents}
\fi
\ifdefined\listfigurename
  \renewcommand*\listfigurename{List of Figures}
\else
  \newcommand\listfigurename{List of Figures}
\fi
\ifdefined\listtablename
  \renewcommand*\listtablename{List of Tables}
\else
  \newcommand\listtablename{List of Tables}
\fi
\ifdefined\figurename
  \renewcommand*\figurename{Figure}
\else
  \newcommand\figurename{Figure}
\fi
\ifdefined\tablename
  \renewcommand*\tablename{Table}
\else
  \newcommand\tablename{Table}
\fi
}
\@ifpackageloaded{float}{}{\usepackage{float}}
\floatstyle{ruled}
\@ifundefined{c@chapter}{\newfloat{codelisting}{h}{lop}}{\newfloat{codelisting}{h}{lop}[chapter]}
\floatname{codelisting}{Listing}
\newcommand*\listoflistings{\listof{codelisting}{List of Listings}}
\makeatother
\makeatletter
\makeatother
\makeatletter
\@ifpackageloaded{caption}{}{\usepackage{caption}}
\@ifpackageloaded{subcaption}{}{\usepackage{subcaption}}
\makeatother
\usepackage{bookmark}
\IfFileExists{xurl.sty}{\usepackage{xurl}}{} % add URL line breaks if available
\urlstyle{same}
% fallback for those not using the hyperref driver hyperxmp:
\makeatletter
\@ifundefined{xmpquote}{\newcommand{\xmpquote}[1]{#1}}{}
\makeatother
\hypersetup{
  pdftitle={What changes when the gas can change density?},
  pdflang={en},
  colorlinks=true,
  linkcolor={blue},
  filecolor={Maroon},
  citecolor={Blue},
  urlcolor={Blue},
  pdfcreator={LaTeX via pandoc}}


\title{What changes when the gas can change density?}
\usepackage{etoolbox}
\makeatletter
\providecommand{\subtitle}[1]{% add subtitle to \maketitle
  \apptocmd{\@title}{\par {\large #1 \par}}{}{}
}
\makeatother
\subtitle{A simple rotating-gas comparison beyond the constant-density
assumption}
\author{}
\date{2026-09-12}
\begin{document}
\maketitle


\href{../paper-map/article.html}{Return to the paper map} · Physical
interpretation, article 3 of 3

The paper studies an incompressible fluid. In its constant-density
setting, pressure adjusts without requiring the density to change.

A real gas responds differently. If its pressure changes substantially,
its density and temperature can change too. What would that do to a
vortex?

To isolate the idea, our companion study compares two descriptions of
the \textbf{same prescribed stationary rotation}: one holds density
fixed, and the other lets an ideal gas respond. We are asking what a
model assumption changes, not simulating the paper's collapse.

\subsection{Begin with the inward push
again}\label{begin-with-the-inward-push-again}

Fluid travelling in a circle needs an inward acceleration. Higher
pressure outside than inside supplies that push.

Faster rotation demands a larger pressure difference. For a fixed
radius, doubling the speed makes the required circular acceleration four
times as large.

In a constant-density model, pressure can adjust while density stays
fixed by assumption. In the gas model, pressure and density are linked
by a stated thermodynamic relationship.

Our comparison uses an isentropic ideal-gas relation. That means the
model links pressure, density, and temperature with a fixed entropy. You
do not need its formula to understand the comparison, but the choice
matters to the numerical results.

\subsection{Hold the rotating shape
fixed}\label{hold-the-rotating-shape-fixed}

We prescribe a smooth circular-speed profile that is zero on the axis,
rises to a maximum, and then decreases outside. Its characteristic
radius is one millimetre.

The outer gas state is fixed at the study's room-temperature reference,
23 °C. We change the maximum speed and calculate the pressure and
density required for stationary radial balance.

Holding the rotation fixed lets us ask a clean question: what difference
comes from allowing density to respond?

It does not tell us how a real vortex would evolve to reach that
prescribed rotation. There is no collapse time or changing core radius
in this experiment.

\subsection{The faster cases have a less dense
centre}\label{the-faster-cases-have-a-less-dense-centre}

Here are three representative results. Density is given relative to the
outer gas:

{\def\LTcaptype{none} % do not increment counter
\begin{longtable}[]{@{}ll@{}}
\toprule\noalign{}
Maximum circular speed & Centre density compared with the outside \\
\midrule\noalign{}
\endhead
\bottomrule\noalign{}
\endlastfoot
About 34.5 metres per second & About 99\% \\
About 103 metres per second & About 88\% \\
About 207 metres per second & About 58\% \\
\end{longtable}
}

The results are specific to the chosen speed profile, outer state, and
gas law. They are not a universal table for every vortex.

\begin{figure}

\centering{

\pandocbounded{\includegraphics[keepaspectratio,alt={Three bars show approximately 99, 88 and 58 percent centre density for maximum speeds about 34.5, 103 and 207 metres per second.}]{figures/reader.png}}

}

\caption{\label{fig-density}The centre-density fraction decreases as the
imposed maximum rotation increases. The values come from the stationary
ideal-gas comparison, not from the time-dependent paper flow.}

\end{figure}%

At the lower speed, constant density is a close approximation for this
measure. At the higher speeds, the same assumption misses a substantial
change.

\subsection{A tolerance makes ``close enough''
explicit}\label{a-tolerance-makes-close-enough-explicit}

Suppose we decide that a five-percent change in centre density is large
enough to deserve attention. That is a chosen tolerance, not a law of
nature.

In this particular example, the five-percent marker occurs at a maximum
speed of about \textbf{67 metres per second}.

Using a tolerance makes the question clearer than asking for one
universal speed at which incompressibility ``breaks.'' A different
velocity shape, state, or required accuracy can give a different answer.

The supporting calculation also compares pressure predictions. They
agree well at low speeds and separate as the gas response becomes
stronger.

\subsection{Why pressure labels
matter}\label{why-pressure-labels-matter}

A pressure below an arbitrary mathematical reference can simply be
negative in that convention. That is common in incompressible
calculations.

Absolute gas pressure has a different interpretation. Once the outer
absolute pressure is fixed, a model that demands an inadmissible inner
absolute pressure has lost its intended physical interpretation.

The main reader-level point is to keep reference pressure and absolute
thermodynamic pressure distinct. A minus sign by itself is not enough to
diagnose a physical failure.

\subsection{Does this show that compressibility prevents
blow-up?}\label{does-this-show-that-compressibility-prevents-blow-up}

This experiment cannot answer that question. It is stationary, has a
fixed radius, and does not evolve the coupled velocity and density in
time.

Its useful finding is narrower: allowing density to change alters the
pressure-supported equilibrium before molecular dimensions are involved.

It also helps separate the two model questions discussed in the
\href{../d9_regime_map/article.html}{previous article}. A flow may still
have useful continuum scale separation while a constant-density
approximation is no longer accurate enough.

The comparison uses a known inviscid equilibrium model. Viscosity and
time-dependent dynamics would require a different experiment.

\subsection{Where this fits in the bigger
picture}\label{where-this-fits-in-the-bigger-picture}

\textbf{A mathematical construction is tied to its equations and
assumptions. Changing the density assumption changes the physical
question.} This comparison makes one such change concrete without
treating it as a critique of the theorem's internal logic.

It completes the short physical-interpretation route. Return to the
\href{../paper-map/article.html}{paper map} to follow the mathematical
mechanism.

The \href{technical-article.html}{technical article},
\href{../../code/d10_compressible_vortex/study.ipynb}{notebook}, and
\href{source-notes.md}{source notes} contain the explicit profile, gas
law, boundary conditions, and independent numerical checks.




\end{document}
